% mono/preamble/eq-tags.tex
%
% Equation-level tag rendering. In segment source these look like:
%
%     *[Derived (slug-name, from prior-slug)]*
%
%     $$ \alpha > \rho/R $$
%
% They mark the epistemic mode of the *next* equation (Derived / Hypothesis
% / Empirical Claim / Formulation / Discussion / Assumption / Postulate /
% Definition). FORMAT.md §Equation-Level Tags lists the full vocabulary.
%
% Rendering: margin note attached to the equation it labels. kaobook's
% \marginnote machinery handles placement, hyphenation, and the
% sans/small typographic register.
%
% Usage from the kramdown emitter:
%
%     \eqtag{Derived (slug, from prior-slug)}
%     \begin{equation*} ... \end{equation*}
%
% The tag attaches to the line where \eqtag is called, which lines up with
% the equation that immediately follows. (kaobook's \marginnote works in
% both horizontal and vertical modes.)

% kaobook's \marginnote (redefined in kao.sty:478) takes the offset as
% its FIRST optional argument: \marginnote[<offset>]{<text>}. kaobook's
% \marginnotevadjust default (0.8\vscale) already pushes all
% marginnotes down a bit globally, and the renderer's lazy-emit logic
% places eq-tags adjacent to their display equation in the source
% stream — so no per-call offset is needed. The chapter-first / subhead-
% direct-to-equation case where the eq-tag floats one line above the
% equation is the cost we accept for not over-shifting the common case.
\NewDocumentCommand{\eqtag}{m}{%
  \marginnote{\footnotesize\sffamily\color{status-soft}#1}%
}
